\section{Alert policies and joint guarantees}\label{sec:theory}
\subsection{Patient-level estimands}
A patient supplies a finite sequence of causally available observations, consent states, and frozen scores $S_1,\ldots,S_T\in[0,1]$. Let $Z$ indicate an incident event with an observed valid warning window $[a,b]$. For a threshold $\tau$, the first-alert policy is
\begin{equation}
 \begin{aligned}
 T_\tau&=\inf\{t:S_t\ge\tau\},\quad R_\tau=\ind\{T_\tau\le T\},\\
 L_\tau&=\ind\{T_\tau\notin[a,b]\},\quad Z=1.
 \end{aligned}
 \label{eq:first}
\end{equation}
An empty set gives $T_\tau=\infty$. Assume $0<\Pr(Z=1)<1$ and $0<\alpha<1$. We target conditional incident risk $r(\tau)=\E[L_\tau\mid Z=1]$ and expected workload $w(\tau)=\E[R_\tau]$, measured in reviews per patient. Their joint constraints are $r(\tau)\le\alpha$ and $w(\tau)\le K$. For $m$ independent patient arrivals in a period, the expected number of reviews is $mw(\tau)$; an expectation bound does not guarantee a finite waiting time.

For comparison, a repeated-alert policy uses every threshold crossing. Its incident missed-event loss is $\ind\{\max_{a\le t\le b}S_t<\tau\}$, which is monotone in $\tau$. A rule that emits only the first crossing has lower review count but a different risk.

\begin{proposition}[First-alert non-monotonicity]\label{prop:nonmono}
Nested threshold crossings do not imply monotone timely first-alert loss. There is a trajectory for which lowering $\tau$ changes the incident loss from zero to one, although the probability of any review is non-decreasing as the threshold decreases.
\end{proposition}
\begin{proof}
Let a pre-window score equal $0.6$, an in-window score equal $0.9$, and all other scores equal zero. At threshold $0.8$, the first alert lies inside the window and loss is zero. At threshold $0.5$, the first alert occurs before the window and consumes the only alert; loss is one. Both thresholds trigger exactly one review. A threshold above $0.9$ yields no alert and loss one. Thus the loss sequence as thresholds increase is $1,0,1$. Any-review events remain nested because they depend only on the overall score maximum.
\end{proof}
This prevents applying the monotone-loss conformal risk-control theorem \citep{crc} to the first-alert loss without additional assumptions. Counting any in-window crossing while reporting only one early review would evaluate a different policy.

\subsection{Population feasibility}
Define the pre-window maximum $A=\max_{t<a}S_t$, with $A=-\infty$ for an empty prefix, the in-window maximum $B=\max_{a\le t\le b}S_t$, and the maximum $C$ over the monitoring interval used to count reviews. For event-positive patients the monitoring interval includes the warning window. Let $\pi=\Pp(Z=1)$ and
\begin{equation}
 \begin{aligned}
 H(\tau)&=\Pp(A<\tau\le B\mid Z=1),\\
 W(\tau)&=\pi\Pp(C\ge\tau\mid Z=1)\\
 &\quad +(1-\pi)\Pp(C\ge\tau\mid Z=0).
 \end{aligned}
 \label{eq:HW}
\end{equation}

\begin{theorem}[Exact feasibility of a first-alert threshold family]\label{thm:feas}
For a fixed threshold domain $\mathcal T$, a jointly feasible first-alert threshold exists if and only if
\[
 \{\tau\in\mathcal T:H(\tau)\ge 1-\alpha,\ W(\tau)\le K\}\ne\varnothing.
\]
If $\mathcal T_\alpha=\{\tau\in\mathcal T:H(\tau)\ge1-\alpha\}$ is nonempty and has a maximum $\tau_\star$, the minimum feasible expected review budget is
\begin{equation}
 K_{\min}=W(\tau_\star).
 \label{eq:kmin}
\end{equation}
Consequently, joint feasibility is equivalent to $K\ge K_{\min}$. Moreover $K_{\min}\ge\pi(1-\alpha)$. If $\mathcal T_\alpha$ is empty, no review budget makes this threshold family feasible.
\end{theorem}
The maximum assumption always holds for a finite nonempty policy grid. For a continuous domain it is a substantive attainment condition; replacing a minimum by an infimum without attainment can make a boundary feasibility assertion false. This characterization is exact but needs the joint law of $A$ and $B$, not merely an hourly ROC curve.

\begin{corollary}[Closed form when early consumption is absent]\label{cor:quantile}
Suppose no alert is possible before the valid window, the positive-patient review interval equals that window, and the conditional distributions of $C$ are continuous, with positive-class CDF $F_1$ strictly increasing at its unique $\alpha$-quantile $q_\alpha$. Let the negative-class CDF be $F_0$. Then
\begin{equation}
 \begin{aligned}
 K_{\min}&=\pi(1-\alpha)+(1-\pi)\{1-F_0(q_\alpha)\},\\
 q_\alpha&=F_1^{-1}(\alpha).
 \end{aligned}
\end{equation}
This is the familiar threshold/ROC workload identity under an explicit no-early-consumption assumption; it is not a general formula for an unrestricted first-alert policy.
\end{corollary}

\subsection{Finite-sample joint certification}
Let $\mathcal T=\{\tau_1,\ldots,\tau_M\}$ be fixed independently of calibration data, for example using training and validation only. Calibration patients are independent and identically distributed from the deployment population, including the consent mechanism. For each policy compute positive-patient losses and all-patient review indicators. Let $n_+$ be the positive calibration count and $n$ the total count. Use exact one-sided Clopper--Pearson binomial bounds \citep{clopper} with tail probability $\eta=\delta/(4M)$ to obtain $(\ell^r_j,u^r_j)$ and $(\ell^w_j,u^w_j)$.

\begin{theorem}[Joint certificate and honest abstention]\label{thm:certificate}
Conditional on any independently trained predictor and frozen grid, with probability at least $1-\delta$ over calibration, all of the following hold simultaneously:
\begin{enumerate}
\item Every $j$ with $u^r_j\le\alpha$ and $u^w_j\le K$ satisfies both population constraints.
\item If every $j$ has $\ell^r_j>\alpha$ or $\ell^w_j>K$, no policy on this grid satisfies both constraints.
\item Any data-dependent selection among certified policies, including the one with smallest workload upper bound, preserves assertion 1.
\end{enumerate}
No monotonicity of $L_\tau$ is required. If neither condition 1 nor condition 2 is met, the output is \emph{undetermined}, not infeasible. For $n_+=0$, the risk interval is $[0,1]$ and no positive-risk claim is certified.
\end{theorem}
The construction applies simultaneous-confidence/LTT reasoning to the first-alert policy in Proposition~\ref{prop:nonmono}. The three outputs distinguish an acceptable operating point, an excluded grid, and a calibration sample that cannot resolve the question.

The certificate is a high-probability statement about the population risk of the selected policy, distinct from the expectation-over-calibration guarantee of ordinary CRC. Model selection performed on calibration would require including every candidate model/policy in the multiplicity family; here model tuning uses validation only. Separate per-model certificates are not automatically simultaneous across models or seeds.

\begin{corollary}[Prespecified group or site cells]
For groups $g=1,\ldots,G$ specified independently of calibration, apply Theorem~\ref{thm:certificate} within each cell using $\delta_g$ with $\sum_g\delta_g\le\delta$. All certified group-conditional risk and workload statements then hold jointly with probability at least $1-\delta$. The groups may overlap. Empty or underpowered cells can remain undetermined. A rule calibrated in one hospital is not certified in another merely by relabeling its population.
\end{corollary}

\subsection{Consent, exchangeability and shift}
Let $U$ denote the complete patient trajectory, $C$ the possibly time-varying consent trajectory, and $V$ any policy randomization. A frozen causal procedure $\Phi(U,C,V)$ produces permitted features, scores, alert decisions, and losses. The procedure may clear cached features and recurrent state after revocation.

\begin{proposition}[Sufficient exchangeability condition]\label{prop:consent}
If the calibration and future tuples $(U,C,V)$ are identically distributed and independent, and the same frozen measurable transformation $\Phi$ is applied to every tuple, their loss and review vectors are identically distributed and independent. Theorem~\ref{thm:certificate} therefore remains valid, including after conditioning on positive-event membership. Consent need not be independent of the outcome.
\end{proposition}
The weaker condition $C\perp Y\mid X$ alone does not guarantee transfer between different consent populations. For a counterexample, let $X=Y\sim\mathrm{Bernoulli}(1/2)$ and score $S=CX$. Calibration consent is always $C=1$, deployment consent always $C=0$, and the threshold is $1/2$. Both populations satisfy $C\perp Y\mid X$, yet positive-patient missed risk changes from zero to one. Conditional ignorability does not make the joint distributions equal.

Weighted conformal methods address particular distribution shifts when suitable likelihood-ratio information is available \citep{weighted}. Appendix~\ref{sec:weighted-proof} gives a bounded-ratio risk/workload certificate derived from Hoeffding's inequality \citep{hoeffding}. Its known-ratio assumption is stronger than fitting a hospital classifier. Our empirical weighted comparison is consequently a diagnostic, not a guaranteed hospital-transfer procedure.
